\begin{abstract}
Mechanical component failures are \emph{grey swans} - rare but physically plausible events that follow known degradation physics yet are critically underrepresented in training data.
Existing supervised prognostic methods require extensive labeled failure trajectories that are expensive and dangerous to collect, limiting their applicability to novel operating conditions.
We introduce a self-supervised framework that learns degradation-aware representations from abundant unlabeled vibration data using two complementary pretraining strategies: (1)~a Joint Embedding Predictive Architecture (JEPA) that captures waveform texture through masked patch prediction, and (2)~temporal contrastive learning that captures spectral centroid shift - the dominant degradation indicator - by exploiting the temporal ordering of run-to-failure episodes.
A hybrid architecture fusing JEPA representations with handcrafted spectral features achieves $\rmse = 0.055 \pm 0.004$ on bearing remaining useful life prediction, a $75.5\%$ reduction versus time-only baselines and $21.4\%$ versus the best handcrafted-only model (5 seeds).
For cross-dataset transfer (FEMTO$\to$XJTU-SY), temporal contrastive pretraining yields $38\%$ improvement over time-only baselines ($p < 0.001$, 10 seeds).
Mechanistic analysis reveals that JEPA and contrastive encoders learn fundamentally different features - waveform periodicity versus spectral dynamics - explaining their complementary strengths.
\end{abstract}
